\section{Conteos finitos y cotas de cuadratura}\label{app:numquad}
\subsection{La comparación con tableros finitos}\label{sec:ajuste-finito}
Los conteos exactos, recogidos en~\cite{oeis} hasta $n=27$~\cite{q27}, permiten observar la
cantidad cuyo límite da el teorema. Para
$10\le n\le27$, la tabla presenta el entero $Q(n)$ y el residuo escalado que resulta al usar el
predictor canónico. Esta última columna depende de un cálculo numérico y se redondea a seis
decimales. El suplemento conserva $\log Q_{\rm pred}(n)$ y el residuo de cada tamaño en el
archivo de datos. La Figura~\ref{fig:counts} muestra los doce últimos valores.

\begin{center}\small
\begin{tabular}{rrr}
$n$ & $Q(n)$ exacto & $n(\log Q(n)-\log Q_{\rm pred}(n))$ [N]\\ \hline
10 & 724 & $0{,}131887$ \\
11 & 2680 & $-1{,}359492$ \\
12 & 14200 & $-0{,}022508$ \\
13 & 73712 & $0{,}356813$ \\
14 & 365596 & $-0{,}988922$ \\
15 & 2279184 & $-0{,}045399$ \\
16 & 14772512 & $0{,}538500$ \\
17 & 95815104 & $0{,}252984$ \\
18 & 666090624 & $0{,}080786$ \\
19 & 4968057848 & $0{,}261961$ \\
20 & 39029188884 & $0{,}415814$ \\
21 & 314666222712 & $0{,}161054$ \\
22 & 2691008701644 & $0{,}101756$ \\
23 & 24233937684440 & $0{,}253993$ \\
24 & 227514171973736 & $0{,}350734$ \\
25 & 2207893435808352 & $0{,}308156$ \\
26 & 22317699616364044 & $0{,}259576$ \\
27 & 234907967154122528 & $0{,}322261$ \\
\end{tabular}
\end{center}
Para ver cuánto de la variación a tamaño finito explica una oscilación de periodo cuatro, el
ajuste archivado usó los tamaños $16\le n\le27$. Fijó $\cerr=0{,}4055004460$, estimado entonces
desde el continuo. Conservamos ese dato para reproducir el ajuste y la Figura~\ref{fig:counts}
con \path{suplemento/datos/ajuste_finito.json}. Difiere en menos de $7\cdot10^{-8}$
del centro actualizado $0{,}40550038$ y cae dentro de la banda conservadora.
Los otros tres parámetros se ajustaron al modelo
\[
 nR_Q(n)=\cerr+\frac{a}{n}+A e^{-\kappa(n-16)}\cos(\pi n/2)+\text{residuo}.
\]
El coseno vale cero para $n$ impar, $+1$ para $n\equiv0\;(\mathrm{mod}\,4)$ y $-1$ para
$n\equiv2\;(\mathrm{mod}\,4)$. El contraste periódico se hace entre las dos clases pares.
El ajuste dio $a=-2{,}956$, $\kappa=0{,}166$ y $A=0{,}297$ [N]. Al introducir la oscilación, la raíz del
error cuadrático medio bajó de $0{,}128$ a $0{,}041$. La corrección $a/n$ con $a$ cercano a
$-2{,}9$ y la oscilación son apreciables para estos tamaños. Los conteos son compatibles con
el valor calculado para $\cerr$, pero la tabla por sí sola no permite determinarlo. El argumento
de dos cortes establece el límite independientemente de este ajuste.

\subsection{Por qué la cota de integración es condicional}
En la Sección~\ref{sec:numeval} se usaron una serie cerca de cero y dos reglas de Gauss en el
resto del intervalo. En los tres tramos hay que distinguir el error del método de integración
del error con que se calculó la función en las muestras. Las cotas de Cauchy controlan los
coeficientes que el muestreo finito mezcla u omite.

En el primer intervalo $[0,r]$, con $r=0{,}03$, escribimos $k_6(x)=\sum k_jx^j$ y $\mathcal V(x)=\sum
v_jx^j$. Los dos coeficientes nulos de $k_6$ dan la integral exacta $\sum_{j\ge2}k_jr^{j-1}/(j-1)$. El
cálculo extrae por separado los coeficientes de Fourier de $d$ y de $\mathcal V$. Si la tilde indica
coeficientes reconstruidos con $N=32$ valores en $|x|=r$, para $2\le j<N$ se tiene
\[
 \widetilde d_j+(1-j)\widetilde v_j-k_j
   =\sum_{\ell\ge1}\bigl(k_{j+\ell N}+\ell N v_{j+\ell N}\bigr)r^{\ell N}
\]
para muestras exactas. En efecto, el muestreo de una función analítica agrega los índices congruentes
con $j$ módulo $N$, y $k_m=d_m+(1-m)v_m$ para $m\ge1$. Cada diferencia se integra con peso
$r^{j-1}/(j-1)$ y la cola de la integral exacta comienza en $j=N$. La fórmula de Cauchy en el círculo
exterior acota los dos tipos de coeficientes. El círculo intermedio de radio $0{,}055$ usa 33 muestras
independientes (64 con las conjugadas) para mejorar las cotas. El radio mayor usado en esta estimación
es $0{,}095$.

En los otros dos intervalos interesa saber qué coeficientes pueden contribuir al error de Gauss.
En un panel real $[a,b]$ sean $h=(b-a)/2$, $y=(a+b)/2$ y $f=k_6/x^2$. Escribimos
\[
 g(w)=f\!\left(y+\tfrac h2(w+w^{-1})\right)
      =\sum_{j\in\mathbb Z}\beta_jw^j .
\]
El cambio $w\mapsto w^{-1}$ deja fijo el argumento de $f$, de modo que $\beta_{-j}=\beta_j$.
Por analiticidad en la elipse de Bernstein de parámetro $\rho_2$ (véase~\cite{trefethen}), la
fórmula de Cauchy da
$|\beta_j|\le\overline M\rho_2^{-j}$ para $j\ge0$, donde $\overline M$ acota la media del
módulo en el contorno. Sea $E_s$ la integral menos la regla gaussiana de $s$ puntos sobre $[-1,1]$.
Para el polinomio de Chebyshev $T_j$ de grado $j$, la simetría y la exactitud de la regla dan
$E_s[T_j]=0$ si $j$ es impar o $j<2s$. Para $j\ge2s$ par, la positividad da
$|E_s[T_j]|\le2+2/(j^2-1)$. El error del panel es, por tanto,
\[
 h\sum_{\substack{j\ge2s\\j\ {\rm par}}}2\beta_j E_s[T_j].
\]
Para estimar estos coeficientes de orden alto se muestrea un contorno interior de parámetro
$1<\rho<\rho_2$ en $N$ raíces de la unidad multiplicadas por $\rho$. La diferencia entre el
coeficiente de Fourier reescalado de las muestras y $\beta_j$ es
\[
 \sum_{\ell\ne0}\beta_{j+\ell N}\rho^{\ell N}+\rho^{-j}\widehat e_j ,
\]
donde $e$ es el vector de errores y la transformada de Fourier se normaliza por $1/N$. El primer
término y la cola omitida se acotan con sumas geométricas en la elipse exterior. Por Parseval,
$\sum_{j=0}^{N-1}|\widehat e_j|^2\le N^{-1}\sum_i\delta_i^2$ si $|e_i|\le\delta_i$.
Para el ruido se retienen los modos $j\le N/2-1$ y el resto de la cola se acota analíticamente. Al
multiplicar por los pesos gaussianos de la expresión anterior, Cauchy--Schwarz acota el ruido. Para el
panel de 20 nodos se emplean 232 muestras en la elipse con $\rho=1{,}64$. Para el de seis nodos, 24
muestras independientes representan una transformada de longitud 46.

Con las tolerancias declaradas de las muestras, las cotas del círculo y de los paneles de seis y veinte
puntos son aproximadamente $7{,}7\,10^{-10}$, $8{,}8\,10^{-11}$ y $1{,}7\,10^{-11}$, respectivamente.
El círculo y el panel corto dan juntos $8{,}56111216\,10^{-10}$. Al sumar el valor archivado del
panel largo, $1{,}69294096\,10^{-11}$, se obtienen unos $8{,}73040626\,10^{-10}$. La cifra comunicada
de $8{,}7\,10^{-10}$ se redondeó al más próximo. Hacia arriba vale $8{,}74\,10^{-10}$ [D si C].
Esta cota sólo cubre la serie y las reglas de integración bajo las tolerancias declaradas.
Para encerrar $\cerr$ hay que propagar además los errores en los 43 nodos del cálculo y en los términos
de frontera. Cuando una tolerancia $\delta_i$ procede de un refinamiento de malla, su valor es una
estimación numérica [N].

